%% !TEX root = manuscript.tex

\chapter{Mission 2 : le Copilote Financier}
\minitoc
\newpage

Ce chapitre présente le projet de Data Science que j'ai conduit au second semestre : le
« Copilote Financier », un outil d'aide à la décision construit sur les données de gestion du
groupe Angelotti, pour sa direction financière. C'est la mission qui se rapproche le plus de ce
que le Master MIASHS m'a préparée à faire, et c'est aussi celle où j'ai le plus mesuré ce que la
structure des données autorise, ou n'autorise pas.

Le périmètre envisagé au départ, la prédiction de la rentabilité des opérations d'aménagement, a
été redéfini en cours de projet. La prédiction de la marge à terminaison subsiste comme premier
axe de travail, mais je l'ai élargie à deux autres questions, et la restitution a finalement été
restreinte à la promotion immobilière. J'expose ici le besoin métier tel qu'il m'a été formulé,
puis les données dont je disposais réellement, l'architecture retenue pour les rendre
exploitables, les trois axes d'analyse et leurs résultats, la plateforme qui les restitue aux
contrôleurs de gestion, et enfin les limites de ce travail.

Une décision a orienté tout le projet. La taille des données disponibles, 267 opérations et
quatre exports, ne justifiait pas un empilement de modèles. J'ai donc retenu un petit nombre de
méthodes simples, chacune choisie pour une question précise, et je me suis attachée à les
valider sérieusement plutôt qu'à les multiplier.

\section{Contexte et besoin métier}

Comme je l'ai indiqué au chapitre~1, chaque opération de promotion est portée par une société
dédiée et vit sur un cycle long. Un budget est arrêté au moment de l'engagement, puis le
chantier et la commercialisation le font dériver au fil des mois : les appels d'offres
reviennent plus chers qu'anticipé, des acquéreurs se désistent, le rythme de vente ralentit sous
l'effet de la conjoncture.

La direction financière suit ces dérives opération par opération, sur tableurs, sans vision
transverse. Personne ne dispose d'une vue d'ensemble permettant de dire, à un instant donné,
quelles opérations méritent l'attention en priorité. C'est ce manque qui a motivé le projet.

\subsection{Les trois questions de gestion}

Trois questions structurent les revues de gestion, et les tableurs actuels n'y répondent pas.

La première porte sur la \textbf{marge} : quelles opérations examiner en priorité parce que leur
marge d'engagement dérive, sur quels postes et pour combien d'euros ?

La deuxième porte sur le \textbf{rythme de vente} : quel volume de réservations budgéter pour
2026 ? La période récente a montré la sensibilité de la demande au coût du crédit, le taux moyen
des crédits habitat étant passé d'environ 1,1~\% à un pic de 3,6~\% fin 2023 avant de refluer
vers 3,1~\% mi-2026, tandis que le rythme du groupe tombait d'environ 150 à moins de 115
réservations par mois en 2023. Reste à savoir ce qui, dans cette chute, revient à la conjoncture
et ce qui revient aux programmes eux-mêmes.

La troisième porte sur le \textbf{prix} : quels lots du stock faut-il réviser, à la hausse ou à
la baisse ? La question n'a rien de théorique, puisque 45~\% des 2~157 désistements enregistrés
ont pour motif un problème de financement ou un refus de prêt, et que les remises sont
aujourd'hui accordées au cas par cas.

\subsection{Besoins par axe et exigences non fonctionnelles}

J'ai formalisé ces besoins en trois axes de travail, complétés d'un axe transverse :

\begin{enumerate}
\def\labelenumi{\arabic{enumi}.}

\item
  \textbf{Axe A, risque de marge} : un score permettant de trier les opérations par risque de
  dérive, le détail en euros des postes déjà en dépassement, et la possibilité de comparer une
  opération à ses semblables.
\item
  \textbf{Axe B, vitesse d'écoulement} : quantifier l'effet du taux d'intérêt sur le rythme de
  réservation, en le séparant de l'effet propre à chaque programme, et simuler des scénarios de
  taux pour 2026.
\item
  \textbf{Axe C, prix} : estimer un prix de marché par lot, signaler le stock qui s'en écarte,
  et proposer une règle de remise argumentée.
\item
  \textbf{Axe transverse} : tirer des signaux simples des commentaires de vente et de la
  structure du réseau de vendeurs.
\end{enumerate}

Quatre exigences non fonctionnelles ont cadré le travail. La \textbf{reproductibilité} d'abord :
tout doit se reconstruire depuis les exports bruts par un seul script, sans étape manuelle
intermédiaire. L'\textbf{interprétabilité} ensuite : chaque score doit être accompagné des
variables qui le déterminent, sans quoi la direction financière ne peut pas le discuter. La
\textbf{confidentialité} également : aucune donnée nominative de client n'est affichée ni
modélisée. Et enfin l'\textbf{acceptation d'un petit échantillon} : à cette taille, je me suis
tenue à des modèles simples, validés en validation croisée.

Le livrable final est une application web destinée à la direction financière, dont la fonction
est de traduire les modèles en décisions. J'en ai restreint le périmètre aux 147 opérations de
promotion immobilière, l'aménagement foncier relevant d'une logique économique différente.

\section{Les données disponibles}

\subsection{Les données internes}

Quatre exports du système de gestion constituent le point de départ, sur un périmètre de 267
opérations.

\begin{longtable}[]{@{}
  >{\raggedright\arraybackslash}p{(\columnwidth - 6\tabcolsep) * \real{0.2523}}
  >{\raggedright\arraybackslash}p{(\columnwidth - 6\tabcolsep) * \real{0.1982}}
  >{\raggedright\arraybackslash}p{(\columnwidth - 6\tabcolsep) * \real{0.2432}}
  >{\raggedright\arraybackslash}p{(\columnwidth - 6\tabcolsep) * \real{0.3063}}@{}}
\toprule\noalign{}
\begin{minipage}[b]{\linewidth}\raggedright
Export
\end{minipage} & \begin{minipage}[b]{\linewidth}\raggedright
Grain
\end{minipage} & \begin{minipage}[b]{\linewidth}\raggedright
Volumétrie
\end{minipage} & \begin{minipage}[b]{\linewidth}\raggedright
Usage
\end{minipage} \\
\midrule\noalign{}
\endhead
\bottomrule\noalign{}
\endlastfoot
Export 1 & lot × dossier de vente & 16~467 lignes, 196 opérations & ventes, prix, dates, désistements \\
Export 2 & poste analytique & 65~000 lignes, 267 opérations & budget, engagé et facturé par poste \\
Export 3 & poste budgétaire & 79~106 lignes & contrôle de cohérence \\
Export 4 & dossier désisté & 2~164 lignes & détail des désistements \\
\end{longtable}

Comprendre ces exports m'a demandé plus de temps que de les modéliser, et c'est normal : c'est à
cette étape que se jouait la fiabilité de tout le reste. Le nettoyage a traité les défauts
classiques d'un export de gestion. J'ai harmonisé les natures de lot en huit familles, fusionné
les doublons orthographiques de communes puisque 91 orthographes désignaient en réalité 90
communes, et converti en valeurs manquantes les surfaces ou prix à zéro, un zéro de surface
n'ayant aucun sens économique.

\subsection{Les données externes}

Trois sources publiques complètent les données internes. Je les ai toutes récupérées par script,
pour qu'elles restent re-téléchargeables et que l'analyse puisse être rejouée sur des séries
actualisées : le \textbf{taux moyen des crédits habitat en France}, série mensuelle de la Banque
centrale européenne~[4] de 2015 à 2026 ; l'\textbf{indicateur de confiance des ménages} publié
par Eurostat~[5] sur la même période ; et un \textbf{référentiel des 90 communes d'implantation}
du groupe, avec population et coordonnées géographiques, obtenu par l'API publique
\texttt{geo.api.gouv.fr}~[6] et enrichi d'une distance au littoral que j'ai calculée.

\section{L'architecture technique}

L'architecture tient en une chaîne courte et entièrement rejouable. Les quatre exports et les
trois sources externes alimentent un script de nettoyage, qui charge une base SQLite organisée
en trois dimensions, les opérations, les communes et la conjoncture mensuelle, et quatre tables
de faits, soit 64~788 lignes budgétaires, 14~123 lots, 14~428 dossiers de vente et 2~164
désistements. Les notebooks d'analyse et l'application s'appuient tous deux sur cette même base.

Trois vues SQL y encapsulent les calculs partagés : la marge budgétée comparée à la marge
engagée par opération, les réservations par opération et par mois, et l'état du stock. Ce choix
présente un avantage pratique : la définition de la marge est écrite une seule fois, dans la
vue, et les notebooks comme la plateforme consomment exactement la même.

Une seule commande reconstruit l'ensemble depuis les exports bruts, en quelques secondes. Tous
les résultats présentés ici sont donc reproductibles, contrainte que je me suis imposée avant
d'écrire la première ligne d'analyse.

\section{Les analyses et leurs résultats}

\subsection{L'exploration et le nettoyage}

L'exploration a produit deux résultats qui ont directement orienté la modélisation.

Le premier concerne la distribution des prix. Le prix au mètre carré des appartements est
fortement asymétrique, avec un coefficient d'asymétrie de 19,2, ce qui le rendait impropre à une
modélisation linéaire directe. Après passage au logarithme, le coefficient tombe à 0,39 et le
diagramme quantile-quantile devient presque rectiligne. La figure~\ref{fig:prixm2} rassemble les
deux distributions. C'est ce constat qui a conduit l'axe C à modéliser le logarithme du prix
plutôt que le prix lui-même.

\begin{figure}[!htbp]
\centering
\includegraphics[width=0.65\textwidth]{fig_3_1_prix_m2_log.png}
\caption{Prix au mètre carré des appartements : distribution brute et après transformation logarithmique}
\label{fig:prixm2}
\end{figure}

Le second est une leçon de méthode, et c'est le passage du projet dont j'ai le plus appris. En
calculant la corrélation entre les réservations mensuelles du groupe et le taux des crédits
habitat, j'obtenais une valeur quasi nulle, de +0,07. Prise au premier degré, elle signifiait que
la conjoncture n'explique rien du rythme de vente, ce qui contredit tout ce que je sais du
marché. En réalité, la série brute est brouillée par la croissance du portefeuille, le nombre
d'opérations en commercialisation ayant fortement augmenté sur la période. Rapportées au nombre
d'opérations actives, les réservations racontent l'inverse : la figure~\ref{fig:intensite} donne
cette série d'intensité corrigée, dont la corrélation avec le taux d'intérêt passe à -0,34 et
celle avec la confiance des ménages à +0,44. Tout l'axe B repose sur cette correction, et sans
elle j'aurais conclu à l'absence d'effet.

\begin{figure}[!htbp]
\centering
\includegraphics[width=0.65\textwidth]{fig_3_2_intensite_vente.png}
\caption{Intensité de vente, mesurée en réservations par opération active : l'effet de la remontée des taux apparaît une fois l'effet de portefeuille neutralisé}
\label{fig:intensite}
\end{figure}

\subsection{La préparation en SQL}

Le livrable de cette étape est la base décrite plus haut : les vues métier partagées, les index,
et un contrôle croisé des agrégats qui retrouve les mêmes totaux par deux moteurs indépendants.
Ce contrôle garantit que la base chargée dit bien la même chose que les exports d'origine, et
c'est le pendant, sur ce projet, des paliers de validation de la mission précédente.

\subsection{Axe A : le risque de dérive de marge}

La difficulté de cet axe n'était pas le modèle mais la définition de la cible. Les exports étant
une photographie à date, sans historique des révisions budgétaires, je ne peux pas observer
comment un budget a évolué, seulement où il en est aujourd'hui.

J'ai donc mesuré la dérive par les dépassements déjà engagés, poste par poste. Les dépenses à
terminaison sont estimées par le maximum de l'engagé et du budget de chaque poste, selon un
principe simple : un dépassement engagé est acquis, et le budget restant sera dépensé. La
variation de marge rapporte ensuite la marge ainsi recalculée à la marge budgétée. J'appelle
cette grandeur la variation de marge « committée », par emprunt à la notion comptable de coût
engagé : elle ne mesure pas une dérive anticipée, mais la part de dérive que les engagements
déjà pris rendent acquise.

Cette définition impose une précaution sur les variables explicatives. Le modèle n'utilise que
le budget initial et les signaux commerciaux, jamais l'engagé ni le facturé, qui définissent la
cible. Sans cette précaution, il aurait « prédit » sa propre définition, ce qui est exactement le
type de fuite de données contre lequel le cours nous met en garde.

Sur le périmètre exploitable, soit les 123 opérations engagées à au moins 60~\% et dont la marge
budgétée dépasse 50~k€, 36,6~\% des opérations n'ont aucun dépassement engagé, 28~\% franchissent
le seuil de dérive matérielle que j'ai fixé à 2~\% de la marge, et deux cas extrêmes dépassent la
moitié de la marge budgétée. La figure~\ref{fig:marge} donne la distribution complète et situe ce
seuil.

\begin{figure}[!htbp]
\centering
\includegraphics[width=0.65\textwidth]{fig_3_3_distribution_variation_marge.png}
\caption{Distribution de la variation de marge « committée » sur les 123 opérations du périmètre, et seuil de dérive matérielle à -2~\%}
\label{fig:marge}
\end{figure}

Je dois signaler ici un résultat que je n'attendais pas. La partie la plus utile de cet axe pour
l'entreprise n'est pas un modèle appris : c'est le tableau descriptif des dépassements engagés
en euros, par opération et par poste, calculé en SQL pur, qui alimente la page « Alertes marge »
de la plateforme. C'est lui que les contrôleurs de gestion ouvrent en premier.

La modélisation vient en complément, comme score de tri. Une régression du taux de variation
atteint un coefficient de détermination de 0,08 sur l'échantillon test, résultat modeste qui
s'explique : la structure du budget initial prédispose à la dérive sans la déterminer. Une
classification des opérations « à risque », dont la dérive dépasse 2~\%, atteint un score F1, qui
combine la capacité à détecter les cas à risque et celle à ne pas en signaler à tort, de 0,30 à
0,34 en validation croisée stratifiée, contre 0 pour la référence majoritaire qui ne détecte
rien. J'ai confronté quatre classifieurs classiques : régression logistique, séparateur à vaste
marge, arbre de décision et forêt aléatoire. À 123 observations, leurs écarts ne sont pas
significatifs ; mon choix s'est porté sur la forêt aléatoire, dont les probabilités alimentent
les niveaux d'alerte.

\subsection{Axe B : la vitesse d'écoulement}

Budgéter les réservations de 2026 suppose de distinguer, dans le rythme de vente, ce qui relève
de la conjoncture et ce qui relève du programme lui-même.

J'ai modélisé le rythme sur un panel opération × mois, soit 3~230 observations portant sur 160
opérations. J'y ai réintroduit les mois à zéro réservation entre la première et la dernière vente
de chaque programme, car les omettre reviendrait à ne compter que les mois où quelque chose s'est
passé, donc à surestimer le rythme.

Un modèle à effets aléatoires sépare alors ce qui revient à l'opération de ce qui revient à la
conjoncture, et deux enseignements en ressortent. D'une part, les deux tiers de la variance
tiennent au programme lui-même, avec un coefficient de corrélation intraclasse de 0,66. D'autre
part, l'effet du taux d'intérêt est fortement significatif : \textbf{une hausse d'un point de taux
réduit les réservations mensuelles d'environ 19~\%}, pour un coefficient de -0,216 associé à une
p-valeur inférieure à 10\textsuperscript{-15}. La régression naïve, qui ignore la structure en
panel, surestime cet effet.

J'ai ensuite voulu vérifier ce résultat par une autre voie. Sur la série agrégée du groupe, qui
compte 114 mois, une régression dynamique à erreurs ARIMA, dont les résidus sont eux-mêmes
modélisés comme une série temporelle, avec sélection par le critère AICc, ne parvient pas à
identifier l'effet du taux : son coefficient y est nul, avec une p-valeur de 0,82. Les variations
du taux sont trop lentes pour être identifiées sur une seule série de cette longueur.

Ce résultat négatif se lit de deux manières. D'abord, il conforte l'idée que la significativité
obtenue sur le panel ne relève pas d'une régression fallacieuse, puisque la voie temporelle, plus
contraignante sur ce point, ne fait ressortir aucun effet mesurable. Ensuite, il montre que les
scénarios doivent s'appuyer sur les deux modèles à la fois. La série temporelle donne la
trajectoire de référence, avec sa dynamique et sa saisonnalité, tandis que l'élasticité estimée
sur le panel convertit chaque scénario de taux en effet sur le rythme. Sur cette base, reportée
en figure~\ref{fig:scenarios}, une détente à 2,5~\% fin 2026 redonnerait de l'ordre de 6~\% de
rythme en moyenne sur douze mois, quand une remontée à 3,5~\% en retirerait 3~\%.

\begin{figure}[!htbp]
\centering
\includegraphics[width=0.65\textwidth]{fig_3_4_scenarios_taux_2026.png}
\caption{Prévision d'intensité de vente à douze mois sous trois scénarios de taux 2026, avec intervalle à 80~\%}
\label{fig:scenarios}
\end{figure}

À l'échelle d'un programme, la courbe des réservations cumulées suit une forme en S, que j'ai
ajustée par une fonction logistique à trois paramètres. Sur les 28 opérations terminées du
périmètre, cet ajustement bat la droite dans 25 cas, comme l'illustre la
figure~\ref{fig:logistiques}. Il fournit surtout un indicateur utile à la trésorerie, le délai
d'atteinte de 90~\% du potentiel de vente, dont la médiane est proche de vingt mois. Sa limite
est connue et je l'ai vérifiée : en début de commercialisation, le potentiel n'est pas
identifiable, et l'outil n'est donc proposé que sur les programmes suffisamment avancés.

\begin{figure}[!htbp]
\centering
\includegraphics[width=0.65\textwidth]{fig_3_5_courbes_logistiques.png}
\caption{Réservations cumulées et courbes logistiques ajustées sur quatre opérations terminées, avec le repère des 90~\% du potentiel}
\label{fig:logistiques}
\end{figure}

\subsection{Axe C : le prix du stock}

Réviser les prix du stock suppose de savoir à quel prix chaque lot devrait se vendre. Le modèle
que j'ai retenu est volontairement le plus simple qui fonctionne : une régression linéaire du
logarithme du prix au mètre carré des 5~064 appartements transigés depuis 2016, expliqué par la
surface, l'étage, l'exposition, le type de produit, l'agence, la commune, la proximité du
littoral et l'année de réservation. Il explique 88,6~\% de la variance sur l'échantillon test,
avec une erreur type d'environ 11~\% sur le prix.

Son intérêt principal n'est d'ailleurs pas sa performance mais sa lisibilité. Ses coefficients se
lisent comme des prix implicites, ce qui permet de les discuter avec la direction commerciale. On
y lit une prime d'étage de +6~\%, une prime d'exposition sud, une décote de 63~\% pour le
logement social qu'expliquent des prix réglementés, et un effet millésime d'environ +19 points
entre 2016 et 2023, suivi d'un plateau. Ce plateau recoupe l'axe B : les prix du neuf n'ont pas
baissé avec la remontée des taux, c'est le volume qui s'est ajusté. Une variante régularisée
n'apporte rien, avec cinq mille observations pour une trentaine de coefficients.

Appliqué aux 127 appartements en stock, ce modèle sert de détecteur : au seuil d'une erreur type,
103 lots sont dans le marché, 8 au-dessus du prix estimé et 16 en dessous. J'ai ensuite posé la
recommandation de remise comme une optimisation sous contrainte : maximiser le revenu attendu à
douze mois, en ajustant chaque prix dans des bornes commerciales de -10 à +5~\%, sous un objectif
d'accélération de l'écoulement, la demande réagissant au prix avec une élasticité fixée à -1.

Le résultat est robuste, et il va à l'encontre de la pratique actuelle. À cette élasticité, la
remise optimale est uniforme en euros sur tous les lots, de l'ordre de 23~k€ par lot sur
l'opération d'application ; elle est donc proportionnellement plus forte sur les petits lots,
alors que la pratique accorde aujourd'hui des remises en pourcentage. Sur l'opération test, qui
compte 18 appartements en stock, viser +8~\% d'écoulement déplace le volume attendu de 6,8 à 7,4
lots pour un revenu quasi inchangé, en recul de 0,3~\%. L'intérêt de la manœuvre n'est donc pas
le revenu immédiat mais la vitesse, et l'économie de frais de portage que l'objectif ne compte
pas.

\subsection{Axe transverse : signaux faibles, texte et réseau}

Trois constats simples, obtenus avec des moyens simples, complètent le dispositif.

Un moteur de recherche par similarité, construit sur les 9~688 commentaires libres de vente,
remonte les dossiers concernés à partir d'une requête en langage courant, comme « refus de prêt
banque ». C'est un outil d'audit des remises non structurées réellement utilisable par le
service.

L'analyse des motifs de désistement par agence révèle ensuite un problème de qualité de saisie
plus qu'un fait commercial : à Salon-de-Provence, 90~\% des motifs sont enregistrés en « Autres
motifs », ce qui rend ces désistements inanalysables. C'est un constat que je n'ai pas cherché,
et qui a donné lieu à une recommandation opérationnelle immédiate.

La structure du réseau de vente, enfin, se révèle très concentrée : deux acteurs, le réseau
interne et un prescripteur externe, portent 51,6~\% des dossiers attribués, avec des taux de
désistement allant de 9~\% pour le premier à 29,6~\% pour le second.

J'ai enfin tenté de prédire le désistement à partir du seul texte du commentaire, avec un
contrôle de fuite soigneux. Cette tentative plafonne à un score F1 de 0,34 : le texte seul ne
suffit pas, et je le traite donc comme un signal d'appoint, non comme un prédicteur.

\section{La plateforme livrée}

L'application est la réponse que j'ai apportée au risque de sur-complexité évoqué en ouverture :
ce qui est livré au métier est simple. Cinq pages, et un périmètre restreint aux 147 opérations
de promotion.

La page \textbf{Vue d'ensemble}, présentée en figure~\ref{fig:vueensemble}, donne la situation en
quelques secondes : 147 opérations suivies, dont 10 en alerte et 2 à surveiller parmi les 68
opérations en cours suffisamment engagées, les 56 autres étant conformes, 421 lots principaux en
stock, un rythme récent de 4,8 réservations par opération et par mois, pour un taux de crédit à
3,11~\%.

\begin{figure}[!htbp]
\centering
\includegraphics[width=0.65\textwidth]{fig_3_6_plateforme_vue_ensemble.png}
\caption{Plateforme, page « Vue d'ensemble »}
\label{fig:vueensemble}
\end{figure}

La page \textbf{Alertes marge}, présentée en figure~\ref{fig:alertes}, classe les opérations par
niveau d'alerte, le score de l'axe A étant traduit en trois niveaux, et affiche pour chacune les
dépassements déjà engagés en euros ainsi que les postes concernés.

\begin{figure}[!htbp]
\centering
\includegraphics[width=0.65\textwidth]{fig_3_7_plateforme_alertes_marge.png}
\caption{Plateforme, page « Alertes marge » : niveaux d'alerte et dépassements engagés en euros}
\label{fig:alertes}
\end{figure}

Les trois autres pages transposent les axes restants. \textbf{Rythme de vente} fait de l'axe B un
simulateur, où un curseur de taux 2026 produit une phrase de résultat fondée sur l'élasticité du
panel. \textbf{Pilotage des prix} affiche le diagnostic du stock et la grille d'ajustement
recommandée. \textbf{Qualité commerciale} reprend les signaux faibles, alerte de saisie comprise.
J'ai vérifié chaque page avec le cadre de test de Streamlit, y compris aux bornes des curseurs et
sur les sélections vides, et l'ensemble se reconstruit depuis les exports bruts.

\section{Difficultés, limites et perspectives}

Ce projet a suivi une logique différente de celle de la mission précédente, en trois temps
séparés : le \textbf{cadrage}, qui formule les questions de gestion avant de regarder les données
et pose les exigences non fonctionnelles ; l'\textbf{exploration et la modélisation}, chaque
question étant traitée par la méthode la plus simple qui y réponde ; et la
\textbf{restitution}, qui traduit les résultats dans le langage de l'utilisateur.

La manière de vérifier mon travail change également. Sur la reprise des données, la question
était de savoir si ce qui était arrivé correspondait à ce qui était parti. Ici, elle est de
savoir si le résultat tient hors de l'échantillon sur lequel il a été construit. J'y ai
systématiquement associé un échantillon test distinct, une comparaison à une référence naïve et
un contrôle des fuites. C'est la comparaison à la référence qui m'a le plus servi, parce qu'elle
seule dit si un score de 0,34 vaut quelque chose.

Les difficultés, quant à elles, relèvent presque toutes de la donnée elle-même.

La principale est structurelle, et elle commande la suite du projet : les données budgétaires
sont une photographie à date, sans historique des révisions. C'est elle qui plafonne l'axe A, et
c'est pourquoi ma première recommandation à l'entreprise est d'historiser les exports mensuels
dans le nouvel ERP. Au bout d'un an, le score statique deviendra un véritable suivi de
trajectoire.

La deuxième est la qualité de saisie, dont le cas de Salon-de-Provence est l'exemple le plus net,
et qui a produit une recommandation opérationnelle immédiate.

La troisième est l'impossibilité d'estimer certains paramètres en interne. C'est le cas de
l'élasticité de la demande au prix, qui a le bon signe dans nos données sans y être
significative, et que j'ai donc dû fixer plutôt qu'estimer.

Les autres limites sont assumées. L'échantillon de l'axe A est petit, 123 opérations, et les
écarts entre modèles y sont indiscernables, ce qui a motivé le choix d'un modèle interprétable
unique. Les intervalles de prévision de l'axe B ignorent l'incertitude sur les taux futurs,
raison pour laquelle la plateforme présente des scénarios plutôt qu'une prévision unique. Le
texte des commentaires, enfin, reste un signal d'appoint. La plateforme couvre par ailleurs la
seule promotion, soit 147 opérations, et exclut l'aménagement foncier, dont la logique
économique justifierait un traitement dédié ; l'analyse, elle, couvre bien les 267 opérations
des exports de gestion.

\section{Conclusion du chapitre}

Le Copilote Financier répond aux trois questions posées en ouverture, mais pas avec la même
force, et il me semble plus honnête de le dire que de présenter un ensemble homogène. L'axe B
est le plus solide : l'effet du taux d'intérêt sur le rythme de réservation est mesuré, séparé
de l'effet propre à chaque programme, et confirmé par un contrôle méthodologique qui aurait pu
l'invalider. L'axe A, à l'inverse, ne dépasse pas le stade du score de tri, et sa partie la plus
utilisée n'est pas un modèle mais une requête SQL.

Ce déséquilibre n'est pas un accident de parcours, et c'est ce que je retiens principalement du
projet. Il découle directement de la donnée disponible : là où j'avais une profondeur
temporelle, le modèle a produit un résultat exploitable ; là où je n'avais qu'une photographie à
date, il n'a pas pu en produire. Les deux recommandations que j'ai formulées, historiser les
exports mensuels et fiabiliser la saisie des motifs de désistement, sont d'ailleurs de même
nature que les chantiers du chapitre~2 : des travaux de structure, faits pour rendre possible
une analyse à venir.
